\documentclass[11pt,a4paper]{article}
\usepackage{curso}
\begin{document}
\portada{Ejercicios --- Lección}{5.2}{p-bits: MTJ y CMOS}
\section{Ejercicios}
\begin{enumerate}[label=\textbf{E\arabic*.}]
\item (B) Con $\tau_0=1$ ns, ¿qué barrera da un tiempo medio de volteo de 10 ns? ¿Y de 10 años?
\item (B) ¿Por qué un pbit necesita $\Delta\sim\kB T$ y una memoria $\Delta\sim40\,\kB T$?
\item (B) ¿Cuántos colores y fases de reloj por barrido necesita una red cuadrada de grado 4?
\item (B) ¿En qué gasta energía un pbit digital que uno magnético no?
\item (I) Un dispositivo de dos estados salta con $k_\pm=\nu\,\sigma(\pm2u)$. Muestre que $P(+1)=\sigma(2u)$ y que el tiempo de correlación es $1/\nu$.
\item (I) Si la barrera tiene una dispersión de $\pm0{,}5\,\kB T$ entre dispositivos, ¿en qué factor varía $\tau$?
\item (A) Simule con Gillespie 1000 pbits con barreras dispersas y mida la distribución de tiempos de correlación.
\end{enumerate}
\section{Soluciones}
\begin{solucion}
\textbf{E1.} $\Delta=\ln10\approx2{,}3\,\kB T$; 10 años $=3{,}2\times10^8$ s: $\Delta=\ln(3{,}2\times10^{17})\approx40{,}3\,\kB T$.

\textbf{E2.} Un pbit debe fluctuar rápido (muchas muestras por segundo); una memoria debe retener años con error mínimo.

\textbf{E3.} 2 colores, 2 fases por barrido.

\textbf{E4.} En generar números pseudoaleatorios y en mover datos a la LUT y al comparador.

\textbf{E5.} Estacionaria $k_+/(k_++k_-)$ y $k_++k_-=\nu$.

\textbf{E6.} $e^{\pm0{,}5}$: entre $0{,}61\tau$ y $1{,}65\tau$.

\textbf{E7.} Distribución log-normal de $\tau$ si la barrera es gaussiana.
\end{solucion}
\end{document}
